\section*{ABOUT ME}
Tiancheng Liu is a computational media researcher and artist working across Creative AI, multimodal interaction, and culturally grounded computing. His work investigates how machines can understand, sense, and participate in human creative practices, with a particular focus on embodied cultural forms such as Chinese calligraphy.
His Ph.D. research is organized around three interrelated layers of creative practice. The first layer asks how machines can analyze and interpret the aesthetic structure of finished artifacts, including layout, form, balance, and stylistic preference. The second layer asks how machines can sense the often invisible processes behind creation, such as gesture, pressure, rhythm, and other implicit micro-actions. The third layer asks how AI systems can join creative practice as collaborators, supporting co-creation, performance, and expressive interaction.
Building upon this framework, he seeks to explore broader possibilities for computing approaches as a medium for interpretation, expression, and collaboration, expanding how humans and machines can sense, make, and imagine together in the cultural and artistic context.
